% report.tex -- compile with: latexmk -lualatex report.tex
% Every figure in the prose is a macro of generated/numbers.tex (tools/status_report.py);
% rerun `python3 tools/status_report.py reports/status-2026-11 --logs <dir of check logs>`.
% report.tex -- compile with: latexmk -lualatex report.tex
\documentclass[11pt,a4paper]{article}

\usepackage[margin=2.2cm]{geometry}
\usepackage{amsmath,amssymb,amsthm}
\usepackage{fontspec}
\usepackage{unicode-math}
\setmainfont{STIX Two Text}
\setmathfont{STIX Two Math}
\setmonofont{DejaVu Sans Mono}[Scale=0.80]
\usepackage{booktabs,longtable,array}
\usepackage{fancyvrb}
\usepackage{luacode}
\usepackage{enumitem}
\usepackage[hidelinks]{hyperref}
\setlist{itemsep=1pt,topsep=3pt}
\setlength{\emergencystretch}{3em}

% ---------------------------------------------------------------- Lean names
% \lean{name}: typeset a Lean identifier verbatim in monospace, with break
% points after '_' and '.', so long names do not overflow table cells.
\begin{luacode*}
-- break after '_' '.' ':' freely; at a lower->upper camelCase boundary with a small penalty
function leanname(s)
  local buf = ""
  local prev = ""
  for _, c in utf8.codes(s) do
    local ch = utf8.char(c)
    if prev:match("^%l$") and ch:match("^%u$") then
      tex.sprint(-2, buf); buf = ""
      tex.sprint("\\penalty100\\hskip0pt\\relax{}")
    end
    buf = buf .. ch
    if ch == "_" or ch == "." or ch == ":" or ch == "/" then
      tex.sprint(-2, buf); buf = ""
      tex.sprint("\\allowbreak{}")
    end
    prev = ch
  end
  tex.sprint(-2, buf)
end
\end{luacode*}
\newcommand{\lean}[1]{\texttt{\directlua{leanname("\luaescapestring{\detokenize{#1}}")}}}
\newcommand{\lbl}[1]{\lean{#1}}   % paper labels (b-prefixed for Barbot)
\newcommand{\dec}[1]{DEC-#1}

% ---------------------------------------------------------------- maths
\newcommand{\R}{\mathbb{R}}
\newcommand{\Z}{\mathbb{Z}}
\newcommand{\N}{\mathbb{N}}
\newcommand{\C}{\mathbb{C}}
\newcommand{\Mink}{\mathbf{M}}
\newcommand{\Dev}{\mathcal{D}}
\newcommand{\Sc}{\widetilde{S}}
\newcommand{\Isom}{\operatorname{Isom}}
\newcommand{\SOo}{\mathrm{SO}_0}
\newcommand{\sing}{\operatorname{sing}}
\newcommand{\reg}{\operatorname{reg}}
\newcommand{\inter}{\operatorname{int}}

\theoremstyle{plain}
\newtheorem{theorem}{Theorem}[section]
\newtheorem{proposition}[theorem]{Proposition}
\theoremstyle{definition}
\newtheorem{definition}[theorem]{Definition}
\newtheorem{reading}[theorem]{Reading}
\theoremstyle{remark}
\newtheorem*{remark}{Remark}


% generated by tools/status_report.py -- do not edit
\newcommand{\numBarbotClosed}{39}
\newcommand{\numBarbotClosedlean}{39}
\newcommand{\numBarbotFindingrows}{32}
\newcommand{\numBarbotFindings}{99}
\newcommand{\numBarbotNlproved}{36}
\newcommand{\numBarbotPartial}{0}
\newcommand{\numBarbotStated}{39}
\newcommand{\numBarbotStatements}{44}
\newcommand{\numBarbotTarget}{39}
\newcommand{\numBarbotTerminal}{5}
\newcommand{\numBenzecriClosed}{35}
\newcommand{\numBenzecriClosedlean}{35}
\newcommand{\numBenzecriFindingrows}{4}
\newcommand{\numBenzecriFindings}{4}
\newcommand{\numBenzecriNlproved}{33}
\newcommand{\numBenzecriPartial}{0}
\newcommand{\numBenzecriStated}{35}
\newcommand{\numBenzecriStatements}{38}
\newcommand{\numBenzecriTarget}{36}
\newcommand{\numBenzecriTerminal}{2}
\newcommand{\numDynisoClosed}{0}
\newcommand{\numDynisoClosedlean}{0}
\newcommand{\numDynisoFindingrows}{0}
\newcommand{\numDynisoFindings}{0}
\newcommand{\numDynisoNlproved}{10}
\newcommand{\numDynisoPartial}{0}
\newcommand{\numDynisoStated}{0}
\newcommand{\numDynisoStatements}{18}
\newcommand{\numDynisoTarget}{18}
\newcommand{\numDynisoTerminal}{0}
\newcommand{\numEgfClosed}{0}
\newcommand{\numEgfClosedlean}{0}
\newcommand{\numEgfFindingrows}{0}
\newcommand{\numEgfFindings}{0}
\newcommand{\numEgfNlproved}{7}
\newcommand{\numEgfPartial}{0}
\newcommand{\numEgfStated}{0}
\newcommand{\numEgfStatements}{16}
\newcommand{\numEgfTarget}{16}
\newcommand{\numEgfTerminal}{0}
\newcommand{\numEttClosed}{57}
\newcommand{\numEttClosedlean}{57}
\newcommand{\numEttFindingrows}{59}
\newcommand{\numEttFindings}{153}
\newcommand{\numEttNlproved}{96}
\newcommand{\numEttPartial}{48}
\newcommand{\numEttStated}{109}
\newcommand{\numEttStatements}{137}
\newcommand{\numEttTarget}{136}
\newcommand{\numEttTerminal}{1}
\newcommand{\numGfnClosed}{94}
\newcommand{\numGfnOpen}{4}
\newcommand{\numGfnPartial}{0}
\newcommand{\numGfnStatements}{98}
\newcommand{\numGxClosed}{143}
\newcommand{\numGxClosedlean}{143}
\newcommand{\numGxFindingrows}{74}
\newcommand{\numGxFindings}{181}
\newcommand{\numGxNlproved}{90}
\newcommand{\numGxPartial}{6}
\newcommand{\numGxStated}{149}
\newcommand{\numGxStatements}{162}
\newcommand{\numGxTarget}{149}
\newcommand{\numGxTerminal}{13}
\newcommand{\numPkgBenzecriDecls}{2\,447}
\newcommand{\numPkgBenzecriFiles}{26}
\newcommand{\numPkgBenzecriLines}{18\,774}
\newcommand{\numPkgBenzecriModules}{27}
\newcommand{\numPkgBenzecriSorryscaffold}{0}
\newcommand{\numPkgBenzecriSorrystrict}{0}
\newcommand{\numPkgCausalDecls}{11\,851}
\newcommand{\numPkgCausalFiles}{194}
\newcommand{\numPkgCausalLines}{127\,174}
\newcommand{\numPkgCausalModules}{195}
\newcommand{\numPkgCausalSorryscaffold}{5}
\newcommand{\numPkgCausalSorrystrict}{0}
\newcommand{\numPkgGxDecls}{11\,887}
\newcommand{\numPkgGxFiles}{45}
\newcommand{\numPkgGxLines}{123\,414}
\newcommand{\numPkgGxModules}{46}
\newcommand{\numPkgGxSorryscaffold}{0}
\newcommand{\numPkgGxSorrystrict}{0}
\newcommand{\numSyracuseClosed}{0}
\newcommand{\numSyracuseClosedlean}{0}
\newcommand{\numSyracuseFindingrows}{0}
\newcommand{\numSyracuseFindings}{0}
\newcommand{\numSyracuseNlproved}{16}
\newcommand{\numSyracusePartial}{0}
\newcommand{\numSyracuseStated}{0}
\newcommand{\numSyracuseStatements}{29}
\newcommand{\numSyracuseTarget}{29}
\newcommand{\numSyracuseTerminal}{0}
\newcommand{\numYamabeClosed}{0}
\newcommand{\numYamabeClosedlean}{0}
\newcommand{\numYamabeFindingrows}{0}
\newcommand{\numYamabeFindings}{0}
\newcommand{\numYamabeNlproved}{16}
\newcommand{\numYamabePartial}{0}
\newcommand{\numYamabeStated}{0}
\newcommand{\numYamabeStatements}{25}
\newcommand{\numYamabeTarget}{25}
\newcommand{\numYamabeTerminal}{0}

\usepackage{xcolor}
\newenvironment{leancode}{\VerbatimEnvironment\begin{Verbatim}[fontsize=\small,frame=single,framesep=2mm]}{\end{Verbatim}}
\newcommand{\unv}{\textsuperscript{\textcolor{red!60!black}{[r]}}}
\newcommand{\oq}[1]{OQ-#1}

\title{The LeanArticles formalizations: status report\\[2pt]
\large Lean libraries, ideas explored and literature formalized}
\author{Technical report on the \texttt{LeanArticles} hub\\
\small host \texttt{node1}; Lean v4.31.0, mathlib \texttt{fabf563a}}
\date{20 November 2026}

\begin{document}
\maketitle

\begin{abstract}
The \texttt{LeanArticles} hub formalizes L.~Brunswic's papers, together with the papers they rest
on, in Lean~4. Between 19 September and 8 November 2026 a series of agent sessions, each under the
author's rulings, built three Lean libraries:
\begin{itemize}
\item \texttt{lean-causal}: topological spacetimes and flat Lorentzian geometry,
  \numPkgCausalDecls{} strict declarations;
\item \texttt{lean-gx}: singular and branched $(G,X)$-manifolds and BTZ geometry,
  \numPkgGxDecls{} strict declarations;
\item \texttt{lean-benzecri}: locally affine and projective manifolds,
  \numPkgBenzecriDecls{} strict declarations.
\end{itemize}
Every strict declaration rests only on Lean's three standard axioms. The strict libraries contain
no \texttt{sorry}, and the only scaffold left holds \numPkgCausalSorryscaffold{} tagged
\texttt{sorry}s.

On these libraries:
\begin{itemize}
\item T.~Barbot's \emph{Globally hyperbolic flat space-times} is closed in full
  (\numBarbotClosed/\numBarbotTarget). Theorems A and B are proved with no hypothesis beyond the
  paper's. The five group-theoretic results the paper cites (Benoist, Bieberbach~I,
  Auslander--Carrière--Dal'bo, Carrière--Dal'bo and Selberg) were proved along the way.
\item Brunswic's \emph{Branched coverings of singular $(G,X)$-manifolds} is closed at
  \numGxClosed/\numGxTarget, including the BTZ uniformization theorem in its ruled form.
\item Brunswic's \emph{Topological spacetimes (I)} stands at \numEttClosed/\numEttTarget.
\item Chapter~1 of Benzécri's thesis stands at \numBenzecriClosed/\numBenzecriTarget.
\end{itemize}
The GFN Bounds library is a frozen external dependency (\numGfnClosed/\numGfnStatements). Four
further papers have ledgers but no Lean yet.

The formalization found many statements false or incomplete as printed. Each is recorded with a
counterexample or a corrected form, and every corrected form was ruled by the author before it
was certified. Exploratory research toward $3+1$ singular spacetimes (the Barbot--Bonsante--Schlenker
review and the \emph{3+1 notes}) is reported separately. It is marked \unv{} wherever it is not
yet proved.
\end{abstract}

\tableofcontents

%======================================================================
\section{Overview}
%======================================================================

\subsection{The hub and its method}

The hub holds two things: a Lean library and the papers rebuilt from it. Each paper goes through
the same four stages.
\begin{enumerate}
\item The paper's statements are extracted into a \emph{ledger} (\texttt{tools/ledger.py}). It
  has one row per \verb|\label|, with the statement's dependencies and its natural-language
  status.
\item Lean declarations are anchored to the rows.
\item A row's \emph{certification status} is computed from the Lean closure, never typed by hand.
\item The paper's story is rewritten fresh in \texttt{STORY.md}. No twelve-word run may be shared
  with the legacy draft.
\end{enumerate}

Four rules hold throughout, and the gates enforce each of them.
\begin{itemize}
\item \textbf{Lean $\ge$ paper.} A declaration narrower than the paper's statement is never its
  anchor. A statement that is false as printed is certified only in a corrected form, and only
  after the author rules that form to be the formalization of record. The form so ruled is
  \emph{marked}; the printed form stays refuted in Lean, usually by an explicit counterexample.
\item \textbf{Strict and scaffold.} Each package has a strict library (\texttt{warningAsError},
  so a \texttt{sorry} fails the build) and a scaffold library, where tagged \texttt{sorry}s are
  allowed. A label counts only if it is closed in the strict library. Every package also runs an
  axiom sweep (\texttt{scripts/AxiomSweep.lean}) over every declaration.
\item \textbf{Buckets.} An open row is classified as one of:
  \begin{itemize}
  \item \texttt{B}: routine;
  \item \texttt{C}: blocked by a named obstruction;
  \item \texttt{D}: deferred, because it needs a layer not yet built;
  \item \texttt{T}: terminal, outside the target, with a reason.
  \end{itemize}
  The closure figure is the number of rows closed in bucket \texttt{A}, over the in-scope rows
  less the terminal ones.
\item \textbf{The model proposes, the author disposes.} Every reading of an ambiguous sentence,
  and every decision to mark a corrected form, is an author's ruling (\texttt{DEC-n} in
  \texttt{registers/DECISIONS.md}). The sessions wrote, proved and audited; each package was
  audited read-only by a session that did not write it before it graduated to strict.
\end{itemize}

\subsection{Closure of every paper of the hub}

Table~\ref{tab:closure} is generated from the ledgers. The GFN Bounds row comes from
\texttt{lean/docs/COVERAGE.md}.

\begin{table}[h]
\centering\small
\begin{tabular}{@{}>{\raggedright\arraybackslash}p{6.2cm}lrrrrr@{}}
\toprule
paper & Lean home & rows & terminal & stated & partial & closure\\
\midrule
\csname @@input\endcsname generated/closure-table 
\bottomrule
\end{tabular}
\caption{Ledger closure per paper (2026-11-20). ``Stated'' counts rows with a Lean anchor;
``partial'' counts those anchored by a disclosed weakening, a corrected form not yet ruled, or a
proper subset of a multi-part statement. Closure is rows closed in bucket A over the
non-terminal rows. For GFN Bounds the four open rows are terminal remarks.}
\label{tab:closure}
\end{table}

\subsection{What changed since the 2 November report}

The previous report (\texttt{reports/barbot-gx/report.pdf}) stood at Barbot 37/39 with five
named externals, gx 108/149 with Section~4 unstated, and no Benzécri. Since then:
\begin{itemize}
\item \textbf{Barbot's theorems became unconditional.} Every cited external was proved:
  \begin{itemize}
  \item Benoist's result, through Gromov's classification on $\mathbb H^d$ (\dec{111});
  \item Bieberbach~I (\dec{110});
  \item Selberg's lemma (\dec{115});
  \item Auslander--Carrière--Dal'bo and Carrière--Dal'bo, for $\Isom(\Mink^n)$ (\dec{121}) and
    then in the printed $\mathrm{Aff}(n,\R)$ generality (\dec{123}).
  \end{itemize}
  This took the closure to \numBarbotClosed/\numBarbotTarget.
\item \textbf{gx's Section~4 was closed.} The brunswic2021cauchy chain (spears, the $C^1$
  surgery, mixed diamonds, absolute maximality of $\reg M$) and the headline
  \lean{GX.btz_uniformization} were proved (\dec{125}, \dec{126}, \dec{129}), as was a full
  singular $(G,X)$-structure on $\widetilde M$ (\dec{130}). \texttt{lean-gx} is entirely
  strict.
\item \textbf{The Benzécri campaign started} (\dec{112}) and brought Chapter~1 to
  \numBenzecriClosed/\numBenzecriTarget.
\item \textbf{The BBS review and the 3+1 research notes} were written (\S\ref{sec:research}).
\end{itemize}

\subsection{Housekeeping done for this report}

Commit \texttt{18d8c16} made the following changes.
\begin{itemize}
\item The ETT scaffold's \lean{minkowski_exists_measure} \texttt{sorry} is closed by the
  identical, already-proved \lean{Causal.Flat.minkowski_frontier_null_measure} (the \dec{65}
  offer).
\item Stale documentation was brought up to date: \texttt{BARBOT.md}, three package READMEs and
  the BBS review.
\item The ledgers were regenerated. The ETT ledger now reads the current Lean.
\item The work register was reconciled with what the commits show was executed, and no ruling
  was taken.
\item Uncommitted, regenerated docs in the read-only \texttt{lean/} submodule were discarded.
\end{itemize}
\S\ref{sec:open} lists what was deliberately \emph{not} done because it needs a ruling.

%======================================================================
\section{The Lean libraries}\label{sec:libs}
%======================================================================

All packages pin Lean v4.31.0 and mathlib \texttt{fabf563a} (\texttt{hub.toml}). They are built
on \texttt{node1} only, and their \texttt{.lake} directories are kept out of Dropbox. The
figures below come from each package's \texttt{make check} on 2026-11-20.

\begin{table}[h]
\centering\small
\begin{tabular}{@{}lrrrrrl@{}}
\toprule
package & strict modules & declarations & lines & strict \texttt{sorry} & scaffold \texttt{sorry} & axioms\\
\midrule
\texttt{lean-causal} (\lean{Causal}) & \numPkgCausalModules & \numPkgCausalDecls & \numPkgCausalLines & \numPkgCausalSorrystrict & \numPkgCausalSorryscaffold & standard\\
\texttt{lean-gx} (\lean{GX}) & \numPkgGxModules & \numPkgGxDecls & \numPkgGxLines & \numPkgGxSorrystrict & \numPkgGxSorryscaffold & standard\\
\texttt{lean-benzecri} (\lean{Benzecri}) & \numPkgBenzecriModules & \numPkgBenzecriDecls & \numPkgBenzecriLines & \numPkgBenzecriSorrystrict & \numPkgBenzecriSorryscaffold & standard\\
\bottomrule
\end{tabular}
\caption{The three libraries. Declarations are the axiom sweep's count, which includes
auxiliary constants. ``Standard'' means \texttt{propext}, \texttt{Classical.choice} and
\texttt{Quot.sound} only. No package declares an \texttt{axiom}.}
\end{table}

\paragraph{Dependencies.}
\begin{itemize}
\item \texttt{lean-causal} depends on mathlib only.
\item \texttt{lean-gx} pins \texttt{lean-causal}. Only \lean{GX.BTZ} and \lean{GX.Lorentz} may
  import \lean{Causal} (\dec{33}, enforced by \texttt{scripts/causal\_import\_gate.py}): gx's
  Section~4 consumes Barbot's classification (\dec{34}).
\item \texttt{lean-benzecri} depends on mathlib only, by design (\dec{112}); a generalization
  layer is planned.
\item Each package is laid out so that moving it into \texttt{lean/} is a \texttt{git mv}. That
  move waits on \oq{1} and \oq{11}.
\end{itemize}

\subsection{\texttt{lean-causal}: topological spacetimes and flat Lorentzian geometry}

\paragraph{The ETT layer} (\lean{Causal.*}) is the theory of \emph{topological} spacetimes:
spaces with a causality, and no metric.
\begin{itemize}
\item Causal relations $J^\pm$, $I^\pm$, $K^\pm$; push-up; strong causality; global
  hyperbolicity; the generated causality and absolute minimality.
\item Causal curves and causal arc-connectedness; admissible (Dieckmann) measures.
\item Morphisms and time functions; gluings; chains; inextensibility; Cauchy surfaces.
\item Geroch's theorem, (i)$\Rightarrow$(ii)$\Rightarrow$(iii) (\lean{geroch_i_imp_ii},
  \lean{geroch_ii_imp_iii}).
\item $(G,X)$-structures; static spacetimes; Minkowski space; the cone models $E_\alpha$;
  singular flat spacetimes under named model hypotheses.
\item \textbf{Levin's theorem.} A continuous time function exists on every second countable
  spacetime with a closed preorder (\lean{Causal.LevinUtility}, proved from mathlib's Urysohn
  machinery). It discharges the model fact \lean{CurvesExtend}.
\item \textbf{Lemma~K.} A compact connected locally Cauchy set of a causally open, second
  countable GH spacetime is a Cauchy surface
  (\lean{isCauchySurface_of_isLocallyCauchy_of_causallyOpen}, \dec{93}).
\item \textbf{Choquet-Bruhat--Geroch.} The universal maximal Cauchy development over any rigid
  singular $(G,X)$-geometry, in Sbierski's order of proof (\lean{Causal.CBG.cbg},
  \lean{cbg_ett}), with the flat instance \lean{Causal.Flat.flat_cbg}.
\end{itemize}

\paragraph{The flat layer} (\lean{Causal.Flat.*}, about two thirds of the package) formalizes
Barbot:
\begin{itemize}
\item Minkowski space $\Mink^n$ and $\Isom(\Mink^n)$; the Penrose boundary and regular domains;
  flat spacetimes as $(\Isom,\Mink)$-manifolds; the developing map; Cauchy-completeness;
  the cosmological time.
\item The nonloxodromic case (\lean{Flat.Nonlox}, 10 modules), the loxodromic case
  (\lean{Flat.Lox}, 26 modules) and the classification (\lean{Flat.Cla}, 19 modules).
\item Five group-theory sub-libraries, each a complete proof of a result Barbot cites
  (\S\ref{sec:lit}): \lean{Flat.Gromov*}, \lean{Flat.Bieberbach*}, \lean{Flat.Selberg*},
  \lean{Flat.ACD*}/\lean{Flat.AffACD*} and \lean{Flat.MaxUnip*}/\lean{Flat.AffUnip*}.
\item The flat Geroch theorem (\lean{flatGerochGH}): every GH flat spacetime carries admissible
  measures.
\end{itemize}

\paragraph{The scaffold} holds \numPkgCausalSorryscaffold{} tagged \texttt{sorry}s, and nothing
in the strict library depends on them:
\begin{itemize}
\item \lean{GerochEquiv}: Geroch (iii)$\Rightarrow$(i), cited by ETT from O'Neill.
\item \lean{SubExtension}, two \texttt{sorry}s: \lean{exists_mem_image_inf}, which feeds
  \lbl{lem:C_ext_connexe}, and \lean{cauchy_pgcd} (\lbl{prop:Cauchy_pgcd}).
\item \lean{CompactCauchy}: Lemma~K in its bare form, without causal openness.
\item \lean{LemmaKSCOpen}: the same with second countability. This is an open generalization the
  paper does not need (\dec{93}).
\end{itemize}

\subsection{\texttt{lean-gx}: singular and branched $(G,X)$-manifolds, BTZ geometry}

\begin{itemize}
\item \textbf{Topology.} Portly and skeletal sets with their dimension theory
  (\lean{GX.Topology.Portly*}); the universal cover; van Kampen (\lean{GX.Topology.VanKampen}).
\item \textbf{Structures.} $(G,X)$-structures and the classical developing theorem
  (\lean{GX.developing_theorem}, \lean{GX.developing_unique}). The bare-map uniqueness the paper
  writes at l.344 is refuted by \lean{developing_unique_fails_for_bare_maps}.
\item \textbf{Spreads and Fox completion} (\lean{GX.Spread.*}); \textbf{singular manifolds},
  almost-everywhere atlases, quotients and suspensions (\lean{GX.Singular.*}).
\item \textbf{Branched coverings} (\lean{GX.Branched.*}): universality, almost-homotopy classes
  of paths, the Galois correspondence (\lean{GX.galois_correspondence}), and the closing
  examples. Antoine's necklace is built explicitly, and its wildness and Antoine's lemma are
  proved through van Kampen.
\item \textbf{Singular II} (\lean{GX.SingularII.*}): unloopable models, the singular developing
  theorem (\lean{GX.singular_developing_theorem}), tame and virtually tame singularities, and
  the spread topology.
\item \textbf{BTZ} (\lean{GX.BTZ.*}, 18 modules):
  \begin{itemize}
  \item the model $E_0$, its rigidity, its causality and the BTZ topology;
  \item CBG for $E_0$-manifolds (\lean{GX.e0_cbg});
  \item holonomy, including a loxodromic element in $2+1$ (\dec{75});
  \item the brunswic2021cauchy chain (spears, surgery, mixed diamonds, absolute maximality);
  \item global hyperbolicity of $\widetilde M$;
  \item the headline \lean{GX.btz_uniformization} and the full structure on $\widetilde M$
    (\lean{GX.btz_manifold_galoisian_maximal_full}).
  \end{itemize}
\item \textbf{The Lorentz bridge} (\lean{GX.Lorentz.Bridge}) identifies gx's
  $(\R^3,\mathrm{null}Q)$ with \lean{Causal.Minkowski 2}.
\item \textbf{Suspensions} (\lean{GX.Suspension.{Conic,Celestial,Affine}}) is the first layer
  of the 3+1 programme (\S\ref{sec:research}).
\end{itemize}

\subsection{\texttt{lean-benzecri}: locally affine and projective manifolds}

\begin{itemize}
\item The projective sphere $S(E)$; its groups $\mathrm{GL},\mathrm{GA},\mathrm{GS},\mathrm{GU}$;
  hemispheres; convexity.
\item Locally projective and locally affine manifolds as charted spaces over a groupoid; jets
  with the étalé topology; quotients; the developing map; the classes of Déf.~14--16.
\item \textbf{Affine cones} (\lean{Benzecri.Cone.*}): the base $V/\R_+^*$, sections, Benzécri's
  cone $F(B)$, and the functorial equivalence
  \lean{coneEquivalence}${}:{}$ \lean{ConeCat (n+1)} $\simeq$ \lean{ProjCat n}.
\item The counterexample module \lean{ConeCex}, for every $n\ge3$.
\item \textbf{Applications}: $B\times S^1$ (\lean{produit_cercle}); products, for two factors
  and for general $p$; the converse; and \textbf{Theorem~A} (\lean{theoremA}).
\end{itemize}

\subsection{\texttt{lean}: GFN Bounds (read-only)}

\texttt{lean/} is the GFNBoundsLean submodule. It formalizes the proofs of Appendices~A, B and H
of \emph{Universality and Convergence of Generative Flows}: universality, the stable bounds, the
Silva et al.\ comparison, and the doubling-graph counterexample.
\begin{itemize}
\item \numGfnClosed{} of \numGfnStatements{} rows are closed. The \numGfnOpen{} open rows are
  terminal remarks.
\item The submodule is pinned at \texttt{2e9865c} (2026-09-19) and is not written from this
  hub (\dec{2}).
\end{itemize}

%======================================================================
\section{Barbot, \emph{Globally hyperbolic flat space-times}}\label{sec:barbot}
%======================================================================

T.~Barbot, J.\ Geom.\ Phys.\ 53 (2005) 123--165, arXiv math/0402257. The labelled copy is
\lean{papers/barbot/legacy/labeled/GHFspace.tex}. The campaign document is
\texttt{lean-causal/docs/BARBOT.md}.

\subsection{Result}

The ledger has \numBarbotStatements{} rows, \numBarbotTerminal{} of them terminal, and
\numBarbotClosed{} of \numBarbotTarget{} are closed. Both classification theorems are proved
with no hypothesis beyond the paper's. In Lean, Barbot's $\Mink^n$ is \lean{Minkowski d} with
$n=d+1$.

\begin{leancode}
theorem barbot_A_unconditional (M : FlatSpacetime.{u} d) (hGH : GH M)
    (hCC : IsCauchyComplete d M) :
    ∃ (M' : FlatSpacetime.{0} d) (p : M' → M), IsFiniteCover d p ∧
      ∃ (d₁ : ℕ) (_ : d₁ ≤ d) (M₁ : FlatSpacetime.{0} d₁) (B : FlatSpacetime.{0} d),
        AMGH (CauchyCompleteClass d₁) M₁ ∧
        (IsTranslationSpacetime M₁ ∨ IsMisnerFamily d₁ M₁ ∨ IsUnipotentFamily d₁ M₁ ∨
          IsCHCentralExtension M₁) ∧
        IsLTPOver B M₁ ∧ ∃ f : M' → B, IsTameEmbedding d f

theorem barbot_B_geroch_unconditional (M : FlatSpacetime.{u} d) (hGH : GH M)
    (hC : CauchyCompact M) (hM : MGH M) :
    ∃ (M₁ : FlatSpacetime.{0} d) (p : M₁ → M), IsFiniteCover d p ∧
      AMGH (CauchyCompleteClass d) M₁ ∧
      (IsTranslationSpacetime M₁ ∨ IsMisnerFamily d M₁ ∨
        (IsCHTorusProduct M₁ ∧ CauchyCompact M₁))
\end{leancode}

\begin{itemize}
\item \textbf{Theorem~A} (\lbl{bthm:caucomp}): a finite cover of a Cauchy-complete GH flat
  spacetime tamely embeds into a linear twisted product over an absolutely maximal model. The
  model is a translation, Misner, unipotent or central-extension Cauchy-hyperbolic spacetime, of
  dimension at most~$n$.
\item \textbf{Theorem~B} (\lbl{bthm:cauclo}): a finite cover of a maximal Cauchy-compact GH flat
  spacetime is absolutely maximal, and is a translation spacetime, a Misner spacetime, or a
  Cauchy-compact Cauchy-hyperbolic $\times$ torus product.
\end{itemize}
Both read Cauchy-completeness through the bridge (CC$'$) of \dec{35}. Theorem~A is under reading
R1 (\dec{43}).

\subsection{The route}

The proof follows the paper's spine:
\begin{enumerate}
\item the developing map and holonomy;
\item the elementary models (§3) and twisted products (§5);
\item the regular domains and cosmological time of §4;
\item the dichotomy between nonloxodromic holonomy (§9) and loxodromic holonomy (§10);
\item the assembly of §11.
\end{enumerate}

The Lean development departs from the paper in four places.
\begin{itemize}
\item \textbf{Torsion.} Barbot's argument silently assumes the holonomy is torsion-free. The
  formalization supplies Selberg's lemma to pass to a finite-index torsion-free subgroup (findings
  F1--F4, \dec{44}). This is why both theorems are stated ``up to a finite cover''.
\item \textbf{Lemma~K.} The sentence ``every closed spacelike hypersurface of a GH spacetime is a
  Cauchy surface'' (Rem.~2.3, §11) is never proved in the paper. It became Lemma~K, proved in
  ETT's setting through Levin's theorem.
\item \textbf{The flat Geroch theorem} is proved (\dec{65}). With it, Theorem~B's Geroch form
  carries no lemma-K data.
\item \textbf{The externals.} Each cited result was first entered as a named \texttt{Prop}
  hypothesis (\dec{36}), never as an axiom, and later proved (\S\ref{sec:lit}). The chain of
  forms is \lean{barbot_A} $\to$ \lean{_noBieb} $\to$ \lean{_min} $\to$ \lean{_unconditional}.
  Each step removes one bundle.
\end{itemize}

\subsection{What the formalization found in the paper}

The ledger carries \numBarbotFindings{} findings on \numBarbotFindingrows{} rows. The ones that
change a statement are:
\begin{itemize}
\item \textbf{Prop.~4.2} is false in the metric reading. It is marked in a two-case form: a volume
  reading for $n\ge3$, a conformal one for $n=2$ (\dec{102}, \dec{104}).
\item \textbf{Prop.~4.10}: a sign typo (\dec{95}).
\item \textbf{Lemma~5.5} is false for $n\ge3$ (\dec{99}).
\item \textbf{Lemma~7.1} is false for $n\ge4$; it is used in its class-by-class form (\dec{98}).
\item \textbf{Cor.~8.4} is false. A unified repair is proved, and the printed form has a
  counterexample (\dec{101}).
\item \textbf{Prop.~8.12} needs $M$ Hausdorff. The counterexample is two copies of $\Mink^n$
  glued along $\{t<0\}$.
\item \textbf{Prop.~9.1} has a gap that needs discreteness (\dec{97}).
\item \textbf{Rem.~2.3} is false as written (\dec{44}).
\item \textbf{The torsion gap} in Theorem~A and in §10.M (F1--F4, the conjugation formula at
  l.3264), fixed with Selberg's lemma.
\item \textbf{Def.~4.4} read literally admits degenerate domains with infinite cosmological time.
  The witness \lean{not_cosmoBoundedAsPrinted} is kept, and the intended \emph{regular} domain
  is used (\dec{88}).
\end{itemize}
None of these defeats the paper's main theorems. All are repaired, and each repair was ruled.

\subsection{What remains}

The \numBarbotTerminal{} terminal rows are off the spine of Theorems A and B, and were not
attempted:
\begin{itemize}
\item \lbl{bthm:convex_cocompact}, which needs Sullivan's structural stability, together with
  \lbl{bdef:geometrically_finite} and \lbl{blem:admissible_cone}, which serve only it;
\item \lbl{bprop:tri}, an explicit thrice-punctured-sphere computation;
\item \lbl{bthm:cmc}, the CMC foliations of §12.
\end{itemize}
\lean{LimitSetMinimal} (Ratcliffe §12) remains a named hypothesis of \lbl{bprop:admlox}, which
Theorems A and B do not reach.

%======================================================================
\section{Brunswic, \emph{Topological spacetimes (I)}}\label{sec:ett}
%======================================================================

\emph{From scratch to Geroch's and Choquet-Bruhat--Geroch's theorems.} The paper builds causal
theory on topological spaces with a causality and no smooth structure, so that singular
spacetimes (cones $E_\alpha$, BTZ lines) are covered as first-class objects.

\subsection{Status}

The ledger has \numEttStatements{} rows. Of these, \numEttStated{} are anchored in Lean and
\numEttClosed{} are closed, so the closure is \numEttClosed/\numEttTarget. A further
\numEttPartial{} rows are partial. Their gaps fall into three groups.
\begin{itemize}
\item \textbf{Proved under named model hypotheses (bucket D).} The singular material of §3--§4
  is proved under \lean{SingularFlatAxioms} or \lean{BTZAxioms}, but the singular instance has
  not yet been supplied from gx. The CBG results (\lbl{theo:CBG} and its §4 lemmas) are
  sorry-free in \lean{Causal.CBG} but are not yet anchored to their rows; \texttt{CBG-DESIGN.md}
  assigns that marking to the ETT campaign.
\item \textbf{Cited or open (bucket C).}
  \begin{itemize}
  \item \lbl{theo:geroch}: (iii)$\Rightarrow$(i) is cited.
  \item \lbl{lem:C_ext_connexe}, \lbl{prop:Cauchy_pgcd}: scaffold \texttt{sorry}s.
  \item \lbl{prop:mesure_dieckmann_Lorentz}: sorry-free since this report's cleanup, but proved
    on Minkowski space only, where the paper states it for strongly causal Lorentzian manifolds.
  \item Three remarks: \lbl{rem:minimalite_non_hereditaire}, \lbl{rem:K_borelien} and
    \lbl{rem:hypotheses_proposition_prop_gx}.
  \end{itemize}
\item \textbf{Readings awaiting rulings.} The open questions \oq{10}, \oq{13}, \oq{14},
  \oq{15}, \oq{16}, \oq{18} and \oq{20} each name the sentence they read.
\end{itemize}
\texttt{STORY.md} is still a skeleton: the words are the author's to write.

\subsection{Findings}

The ledger carries \numEttFindings{} findings on \numEttFindingrows{} rows. The notable ones:
\begin{itemize}
\item \lean{BTZAxioms.Iplus_eq} (ETT.tex l.1351) is false in $E_0$. It is replaced by
  \lean{Iplus_of_lt} (\dec{64}).
\item The order of quantifiers in Cauchy surfaces and Geroch's (ii) (\dec{14}, \dec{15}).
\item Maximal chains do not exist in Minkowski space as defined; \lbl{def:maximal_chain} is
  superseded by causal completeness through inextensible curves (\dec{20}--\dec{22}).
\item The Baire hypothesis of \lbl{theo:causalite_engendre} is not used (\oq{10}).
\end{itemize}

%======================================================================
\section{Brunswic, \emph{Branched coverings of singular $(G,X)$-manifolds}}\label{sec:gx}
%======================================================================

arXiv 2010.10610, Geom.\ Dedicata 2024. The ledger reads the accepted version of 2023
(\dec{24}); the 2021 pre-referee draft is kept in \texttt{legacy-2021/}.

\subsection{Status}

The ledger has \numGxStatements{} rows, \numGxTerminal{} of them terminal: these come after
\verb|\end{document}| (weak morphisms, \dec{24}). Of the rest, \numGxClosed{} of
\numGxTarget{} are closed, and the \numGxPartial{} remaining rows are partial \emph{by ruling}:
\begin{itemize}
\item \lbl{ex:following_example_skeletal_itemize}: item~3 is false (the bow-tie); its corrected
  form is proved (\dec{27}).
\item \lbl{theo:Fox_completion_functor}: proved on a locally connected Hausdorff base. Existence
  and full faithfulness are false as printed (\dec{48}).
\item \lbl{prop:spread_fiber}: ``injective open'' is read as an embedding; the open-map reading
  is false (a comb, \dec{48}).
\item \lbl{prop:branched_covering_composition}: needs finite fibres or slsc. The counterexample
  is on $D_\infty$ (\dec{48}).
\item \lbl{rem:definition_semi_locally_simply_connected}: the ``iff'' is false (the harmonic
  archipelago). Whether a compact counterexample is wanted is an open question.
\item \lbl{cex:start_singular_manifold_homeomorphic}: Troyanov's theorem is a named external
  (\lean{GX.TroyanovSuspensionTheorem}) with non-vacuity witnesses (\dec{69}).
\end{itemize}

\subsection{The headline: uniformization of BTZ-manifolds}

\lbl{theo:globally_hyperbolic_cauchy_complete_cauchy_maximal} closes in the form ruled by
\dec{28}, \dec{107}, \dec{119} and \dec{129}. Let $M$ be a GH, Cauchy-maximal $E_0$-manifold,
and let $(D,\rho)$ be its developing pair. Suppose:
\begin{itemize}
\item $M$ is Cauchy-\emph{compact}, in the $C^1$ sense;
\item $\pi_1$ of the regular part is free and non-abelian.
\end{itemize}
Then:
\begin{itemize}
\item $D$ is injective, and $\rho$ is faithful, discrete and torsion-free;
\item the image of $D$ is $(\bigcup_i\Delta_i)\cup\bigcap_jI^+(\Pi_j)$;
\item $M\cong\rho\backslash D(\widetilde M)$ as an almost-everywhere isomorphism.
\end{itemize}
The Lean statement is \lean{GX.btz_uniformization} (\texttt{GX/BTZ/Uniformization.lean}). The
paper's ``Cauchy-complete'' is recorded as a finding: Barbot's equality needs absolute
maximality, which brunswic2021cauchy provides only in the Cauchy-compact case.

The route runs through the regular part:
\begin{enumerate}
\item $\reg M$ is flat, GH and absolutely maximal, by brunswic2021cauchy Prop.~2.16--Cor.~2.18,
  proved through spears, the $C^1$ surgery and mixed diamonds.
\item Barbot's Theorem~A, imported from \texttt{lean-causal}, then classifies $\reg M$.
\item The BTZ lines are added back through the BTZ half-line package.
\end{enumerate}
In the natural topology the lifted $\widetilde M$ is globally hyperbolic. With $T_{\mathrm{BTZ}}$
it is not: \lean{btzJ0_not_isCompact_diamond} (\dec{106}, reversing \dec{77}b).

\subsection{Findings against the paper}

The ledger carries \numGxFindings{} findings on \numGxFindingrows{} rows. Those that change a
statement:
\begin{itemize}
\item \lbl{ex:conical}: the printed metric has cone angle $\sqrt\alpha$, not $\alpha$
  (\dec{47}).
\item \lbl{theo:ramified_covering_path} and \lbl{theo:Galois_cor} need the added hypothesis
  ``$Y\setminus S$ semilocally simply connected'' (\dec{70}).
\item \lbl{theo:local_virt_tame} is false in the paper's generality. The counterexample is
  \lean{GX.ThinCex}, a cone over $T^2$. The row is certified in the descending-chart form
  (\dec{109}).
\item \lbl{ex:consider_denote_absolute_galois}: the absolute Galois group is $\Gamma_2'$, the
  kernel of the windings mod~2 (\dec{69}).
\item \lbl{cex:proposition_prop_branched_covering} must use $S\cup\{\text{apex}\}$
  (\dec{53}).
\item \lbl{ex:quotient_base}: the tachyon and photon quotients under the radial action are not
  Hausdorff. The repair is \lean{GX.ThetaLiftQuotient} (\dec{116}).
\item \lean{E0TypeAgreement} is false for $\Mink\oplus E_0$.
\item brunswic2021cauchy: the gluing of Lemma~2.3 is proved globally hyperbolic (\dec{108}). The
  paper's citation of ``Lemma~2.4'' points at the wrong lemma.
\end{itemize}
Proved here rather than cited:
\begin{itemize}
\item the loxodromic element of the holonomy in $2+1$ (\dec{75});
\item Antoine's lemma and the wildness of Antoine's necklace (\dec{82}, \dec{85});
\item the mixed diamonds of Prop.~2.17 (\dec{126}).
\end{itemize}

%======================================================================
\section{Benzécri, \emph{Sur les variétés localement affines et localement projectives}}\label{sec:benzecri}
%======================================================================

J.-P.~Benzécri's thesis, Bull.\ SMF 88 (1960) 229--332. No \TeX{} source exists. The ledger reads a
transcription of the introduction's results and Chapter~1, made on 2026-11-03 from the Numdam scan
(\texttt{papers/benzecri/legacy/labeled/benzecri1960\_ch1.tex}). The author has still to check the
transcription. The campaign is the seed of the generalization layer that the BBS programme needs
(\S\ref{sec:research}, N7).

\subsection{Status}

The ledger has \numBenzecriStatements{} rows, \numBenzecriTerminal{} of them terminal
(results~B and~C belong to Chapters~2 and~5, which are not transcribed). \numBenzecriClosed{}
of \numBenzecriTarget{} are closed.
\begin{itemize}
\item \textbf{Theorem~A} (\lean{theoremA}). If $V$ is a compact Hausdorff locally projective
  $n$-manifold, then $V\times S^1$ is locally affine. Moreover, a product of $q+1$ such manifolds
  with a $q$-torus is locally projective of dimension $\sum n_i+q$.
\item \textbf{The \S2 functorial isomorphism} is \lean{coneEquivalence}, an equivalence of
  categories between affine conical manifolds and locally projective manifolds.
\item \textbf{The one open row} is \lbl{bzthm:c1s1_revetement_universel_divisible} (p.~238): the
  deck group of a compact $V$'s universal cover acts properly discontinuously and cocompactly.
  It needs the universal cover, which is not in mathlib at the pin. \texttt{lean-gx} builds one,
  but \texttt{lean-benzecri} is mathlib-only until the generalization layer exists.
\end{itemize}

\subsection{Findings}

\begin{itemize}
\item \textbf{The base of an affine cone need not be Hausdorff.} The counterexample is
  $V=\{z>0,(x,y)\ne0\}\times\R^{n-3}/\langle\mathrm{diag}(4,1/8,2,1,\dots)\rangle$, for every
  $n\ge3$. The base is Hausdorff iff the $\R_+^*$ action is proper, and always for $n\le2$.
  §2 Prop.~2, Déf.~3 and the functorial isomorphism are marked with that hypothesis added
  (\dec{122}, \dec{124}, \dec{127}).
\item l.640 has a misprint: ``$\varphi\circ j_2=j_1$'' should be $j_2\circ\varphi=j_1$
  (\dec{127}).
\item \textbf{§3 Prop.~2, proof remark.} The product-one lift is not conical when the factors'
  dimensions differ. The corrected lift is $\prod\lambda_i^{\dim F(B_i)}=1$; the statement itself
  holds (\dec{128}).
\end{itemize}

%======================================================================
\section{The other papers of the hub}\label{sec:others}
%======================================================================

The papers below have a frozen legacy snapshot and a generated ledger with its DAG. None has Lean
yet, and their \texttt{STORY.md} headlines are the author's to write.

\begin{itemize}
\item \textbf{EGF}, \emph{Ergodic Generative Flows} (ICML 2025). \numEgfStatements{} rows,
  \numEgfNlproved{} of them proved or definitional in the paper. \texttt{STORY.md} names the
  anchor \lean{EGF.main_universality}, which does not exist yet. EGF is the author's first
  priority (\dec{4}).
\item \textbf{Syracuse}, \emph{On non-trivial cycles of the $3x+1$ problem}: an extension of
  Terras's method. \numSyracuseStatements{} rows.
\item \textbf{dyniso}, \emph{Orthogonal hyper-connections}: why the Birkhoff polytope breaks
  dynamical isometry. \numDynisoStatements{} rows.
\item \textbf{Yamabe}, \emph{The Yamabe gap}: inhomogeneity and cosmic expansion.
  \numYamabeStatements{} rows.
\item \textbf{nonacyclic}, \emph{Non-acyclic GFlowNets}. Blocked: the \TeX{} source has not been
  located (\oq{3}). Only a PDF and a \texttt{ledger.toml} exist.
\end{itemize}

%======================================================================
\section{Ideas explored and ongoing research}\label{sec:research}
%======================================================================

This section reports work that is \emph{not} certified. Claims marked \unv{} are a session's own
reasoning, which nobody has proved.

\subsection{The BBS review and a roadmap to $3+1$}

\lean{papers/bbs/REVIEW.md} (2 November) reads Barbot--Bonsante--Schlenker, \emph{Collisions of
particles in locally AdS spacetimes~I} (arXiv 1010.3602), against the gx library.

In BBS, a neighbourhood of a singular point is the suspension of its link, a singular
$\mathrm{HS}^2$-sphere. The review proposes a layered architecture:
\begin{itemize}
\item \textbf{L0}: a general theory of suspensions;
\item \textbf{L1}: links in locally affine spaces;
\item \textbf{L2}: links by signature;
\item \textbf{L3}: a Riemannian branch and a Lorentzian branch, the latter inheriting ETT's
  causal structure.
\end{itemize}
It also proposes waves W0--W5 and questions BBS-Q1--Q8.

The author's rulings so far:
\begin{itemize}
\item flat first;
\item the full $3+1$ stratification is in scope, in codimensions 2, 3 and 4;
\item stabilisers are pointwise;
\item the Lorentzian suspension is a cocycle on \lean{Susp L}.
\end{itemize}

\textbf{Lean so far.} L0/L1 are built as \lean{GX.Suspension.{Conic,Celestial,Affine}}, audited
twice and graduated. Two results there shaped the definitions:
\begin{itemize}
\item \lean{GX.celestial_counterexample}: the suspension proposition of the thesis draft's
  commented-out material fails when the stabiliser does not preserve the radius. This is a design
  motivation, not a finding against the thesis.
\item \lean{GX.not_conicStructure_of_contraction}: radius preservation is incompatible with
  boosts.
\end{itemize}
A codimension-2 interface for 3+1 would change the ETT and gx interfaces, and needs the author's
ruling (BBS-Q7).

\subsection{The 3+1 notes}

\lean{papers/bbs/NOTES-3+1.md} is an append-only log of entries N1--N8 (2--3 November). In each
entry the author proposes an intuition and the session answers with what the flat $(G,X)$ setting
forces.
\begin{itemize}
\item \textbf{N1: knotted string traces.} The manifold stays standard; the pair $(M,S)$ is what a
  knot spoils. A minimal surface could replace the cone\unv.
\item \textbf{N2: two competing theories} (conical singularities versus a Ricci-flat sector),
  physical selection principles, and a conjectured bridge\unv. Eight literature references are
  not yet archived.
\item \textbf{N3--N4, Theorem~M} (momentum and energy conservation at a $2+1$ interaction)\unv.
  The holonomy of the incoming circle $C_-$ is the ordered product of the particles' momenta.
  Total energy is conserved \emph{as a real number} on the universal cover $\widetilde G$, not
  just mod $2\pi$. Positivity of the total mass is BBS's Def.~4.2. N4 also gives a $2\to2$
  outcome law.
\item \textbf{N5, Theorem~G}\unv. In a positive causal interaction, the type of the total
  incoming momentum decides the outcome:
  \begin{itemize}
  \item timelike: scattering is possible;
  \item spacelike (a Gott pair): the system collapses to a single BTZ-like singularity, a black
    hole;
  \item parabolic: the extreme case.
  \end{itemize}
  This rests on a re-hypothesised BBS Prop.~\texttt{pro.hyphyp} and on a ``funnel lemma'' that is
  still to be written. N5 also gives a $2\to3$ outcome law.
\item \textbf{N6--N7}: what Seppi's work gives the cocycle suspension, and Benzécri's cones as
  the systematic framework for suspensions. N7 motivated the Benzécri campaign.
\item \textbf{N8}: a correction from gx: radial cocycles and Hausdorff quotients.
\end{itemize}
The next steps the notes name are:
\begin{itemize}
\item write out Theorem~M in $2+1$ from BBS §5.1/§7.2, then formalize it in gx (holonomy of a
  punctured $\mathrm{HS}$-sphere);
\item write out the funnel lemma;
\item archive the cited physics literature (Gott; Carroll--Farhi--Guth; Deser--Jackiw--'t~Hooft;
  Matschull).
\end{itemize}
\lean{papers/bbs/SEPPI-READING.md} holds quoted reading notes on seven of Seppi's papers.

\subsection{Other directions}

\begin{itemize}
\item \textbf{The double DAG} (\texttt{reports/doubledag/gfnbounds/}) compares the statement DAG
  and the Lean DAG of GFN Bounds edge by edge. Projecting through \emph{principals} removes 55\%
  of the Lean-only edges and keeps 66\% of the agreements. It is not yet run on the geometry
  papers.
\item \textbf{Shrinking Troyanov's external}: a plan in
  \lean{lean-gx/campaign/POSTWAVE5/BBS-NOTES.md}, waiting for the author.
\end{itemize}

%======================================================================
\section{Literature formalized}\label{sec:lit}
%======================================================================

Table~\ref{tab:lit} lists external results that were proved in full in the strict libraries, and
says where. ``Proved'' means sorry-free with standard axioms.

\begin{longtable}{@{}>{\raggedright\arraybackslash}p{4.6cm}>{\raggedright\arraybackslash}p{4.4cm}>{\raggedright\arraybackslash}p{5.6cm}@{}}
\caption{External results formalized.}\label{tab:lit}\\
\toprule
result & reference & Lean\\
\midrule\endfirsthead
\toprule result & reference & Lean\\ \midrule\endhead
Auslander's theorem, Carrière--Dal'bo form (Barbot Thm~5.1), printed $\mathrm{Aff}(n,\R)$ generality
  & Carrière--Dal'bo; Auslander; via Zassenhaus and Hall's lemma
  & \lean{affAuslanderCarriereDalbo}, \lean{auslanderCarriereDalbo} (\dec{121}, \dec{123})\\
Maximal unipotent subspace (Barbot Prop.~5.2), any field
  & Fried--Goldman--Hirsch (Fitting submodule); Hirsch's $H^1(N;V)=0$ for nilpotent $N$
  & \lean{affMaxUnipotent}, \lean{maxUnipotentSubspace}\\
Selberg's lemma, every finite-dimensional real $V$
  & Selberg; congruence subgroups, $\Z$ Jacobson (proved)
  & \lean{selbergLemma} (\dec{115})\\
Bieberbach's first theorem
  & Bieberbach; virtual abelianity via Zassenhaus
  & \lean{bieberbachI} (\dec{110})\\
Discrete subgroups of $\mathrm{SO}_0(1,d)$ without loxodromics fix a timelike or null vector (in place of Benoist)
  & Gromov's classification of isometries of $\mathbb H^d$
  & \lean{benoistNoLox}, \lean{bprop_sanloxo_proved} (\dec{111})\\
Continuous utility for closed preorders
  & Levin
  & \lean{Causal.LevinUtility}\\
Geroch's splitting, flat case; (i)$\Rightarrow$(ii)$\Rightarrow$(iii) in general
  & Geroch
  & \lean{flatGerochGH}, \lean{geroch_i_imp_ii}, \lean{geroch_ii_imp_iii}\\
Universal maximal Cauchy development
  & Choquet-Bruhat--Geroch, in Sbierski's form (arXiv 1309.7591)
  & \lean{Causal.CBG.cbg}, \lean{flat_cbg}, \lean{GX.e0_cbg}\\
Classical developing theorem for $(G,X)$-manifolds
  & Thurston; Goldman
  & \lean{GX.developing_theorem}\\
Fox completion of spreads (corrected)
  & Fox 1957
  & \lean{GX.Spread.Completion}\\
Antoine's lemma; wildness of Antoine's necklace
  & Antoine; van Kampen
  & \lean{GX.Branched.GaloisExamples}, \lean{GX.Topology.VanKampen}\\
brunswic2021cauchy, Lemmas 2.3--2.18
  & Brunswic, arXiv 2108.12631
  & \lean{GX.BTZ.{Spear,Surgery,MixedDiamond,AbsMax}}\\
Barbot, Theorems A and B
  & Barbot 2005
  & \lean{barbot_A_unconditional}, \lean{barbot_B_geroch_unconditional}\\
Benzécri, Chapter~1 and Theorem~A
  & Benzécri 1960
  & \lean{Benzecri.theoremA}, \lean{coneEquivalence}\\
\bottomrule
\end{longtable}

\paragraph{Still cited, not proved.}
\begin{itemize}
\item Troyanov's theorem on cone metrics (\lean{GX.TroyanovSuspensionTheorem}, a named
  hypothesis, \dec{69}).
\item Geroch (iii)$\Rightarrow$(i), from O'Neill (an ETT scaffold \texttt{sorry}).
\item \lean{LimitSetMinimal}, from Ratcliffe §12 (a named hypothesis off Barbot's spine).
\item Sullivan's structural stability, for Barbot's terminal \lbl{bthm:convex_cocompact}.
\end{itemize}

\paragraph{Read but not formalized.}
\begin{itemize}
\item BBS~I (arXiv 1010.3602) and BBS~II (1202.5753);
\item seven papers by Seppi (\lean{papers/bbs/sources/seppi/});
\item arXiv 1605.05530 (BTZ extensions);
\item Brunswic's thesis.
\end{itemize}

%======================================================================
\section{Open items for the author}\label{sec:open}
%======================================================================

These items need a ruling or the author's own words. None was acted on here.

\begin{enumerate}
\item \textbf{The freeze has expired, but is still in place.} \texttt{hub.toml [freeze].until}
  is 2026-09-25. \texttt{CLAUDE.md}, \dec{2}, and risk R-1 (the ICLR deadline) all still describe
  it. Consequently \texttt{tests/test\_no\_model\_in\_gates.py}, which runs the gate with
  \texttt{FREEZE=1}, now fails \texttt{gate-registers}, so \textbf{the hub's \texttt{make check}
  is red}. Retiring the freeze, together with \oq{1} (absorbing GFNBoundsLean in-tree), is a
  ruling.
\item \textbf{\lean{papers/bbs} fails \texttt{papers\_check}} because it has no ledger. It needs
  either a ledger over BBS~I or a ruling that a review directory is not a paper. This, too, keeps
  \texttt{make check} red. The per-package gates of \texttt{lean-causal}, \texttt{lean-gx} and
  \texttt{lean-benzecri} are green.
\item \textbf{\texttt{hub.toml [papers].order}} (\dec{4}) omits \texttt{barbot},
  \texttt{benzecri} and \texttt{bbs}. They are appended silently.
\item \textbf{Work packages awaiting review.} Most of the work from 26 September onward is
  \texttt{IN-REVIEW} in \texttt{registers/WORK.md}:
  \begin{itemize}
  \item ETT: WP-15, 18, 20--22;
  \item gx: WP-24--28, 38--41, 43, 44, 46--48;
  \item Barbot: WP-29--34;
  \item Benzécri: WP-45.
  \end{itemize}
  WP-19 and WP-36 are still \texttt{IN-PROGRESS}. \texttt{registers/QA.md} has no rows: no
  evaluation has been recorded.
\item \textbf{ETT open questions} \oq{10}, \oq{11}, \oq{13}--\oq{16}, \oq{18} and \oq{20}.
  \oq{11} asks whether \texttt{lean-causal} remains a package of its own.
\item \textbf{The CBG marking.} Anchoring \lean{Causal.CBG} to \lbl{theo:CBG} and its §4 lemmas
  would lift a block of ETT rows.
\item \textbf{The Benzécri transcription} has to be checked against the scan before its
  certifications are cited.
\item \textbf{The \texttt{STORY.md} headlines} of ETT, gx, Syracuse, dyniso and Yamabe are
  skeletons. EGF has no Lean.
\item \textbf{nonacyclic}: the source has to be located (\oq{3}).
\item \textbf{The 3+1 programme.} The open questions BBS-Q1--Q8 and N1-Q, N2-Q, together with
  Theorem~M's formalization, are the natural next campaign. BBS-Q7 (codimension-2 strata) is an
  interface ruling.
\end{enumerate}

\paragraph{Provenance.} The figures come from \texttt{generated/numbers.tex}, written by
\texttt{tools/status\_report.py} from the ledgers, \texttt{lean/docs/COVERAGE.md}, each package's
\texttt{scripts/sorry\_count.py}, and the axiom sweeps of the \texttt{make -C <pkg> check} runs of
2026-11-20 on \texttt{node1}. Rulings are quoted by number from \texttt{registers/DECISIONS.md}.

\end{document}
